\begin{center}
\textcolor{minesblue}{\rule{0.6\textwidth}{1pt}}\\[4pt]
{\Large\bfseries\sffamily\color{minesblue} Phase 3: ``How Well Did It Work?''}\\[2pt]
{\small\itshape Sprints 5--7 execute V\&V, collect evidence, and deliver the final package.}\\[4pt]
\textcolor{minesblue}{\rule{0.6\textwidth}{1pt}}
\end{center}
\vspace{0.5em}
%=============================================================================
\section{Sprint 5 -- Complete System Demonstration}
%=============================================================================

\begin{priorities}
\begin{enumerate}[leftmargin=*, itemsep=1pt, topsep=0pt]
    \item Demonstrate end-to-end functionality from the user's perspective.
    \item Update the BoM with actual costs and a labor estimate.
    \item Finalize the test plan so it is ready for the external team---this is your last chance to fix ambiguities.
\end{enumerate}
\end{priorities}

\begin{faqbox}[Q: What should I focus on for the in-class system demo?]
Show that your system works \textit{end-to-end} as a user would experience it. Start with the user powering on the system, walk through the primary use case, and demonstrate key requirements being met. If something is not working, be upfront about it. Instructors value honesty and a clear understanding of what went wrong far more than a polished demo that hides problems.
\end{faqbox}

\begin{faqbox}[Q: My design has changed a lot since the CDR. Is that a problem?]
Not at all---design iteration is expected and encouraged. What matters is that you \textit{document} the changes: what changed, when it changed, and why. Part A of the Complete System Demo Report specifically asks for this. A well-documented design evolution shows engineering maturity.
\end{faqbox}

\begin{protip}
\textbf{When debugging integration failures, follow the five-step protocol from Chapter~20 of the MRD.} (1)~\textbf{Reproduce}: make the failure happen reliably and note the exact conditions. A failure you cannot reproduce is a failure you cannot fix. (2)~\textbf{Isolate}: disconnect modules one at a time until the failure stops---the last module disconnected is the prime suspect. (3)~\textbf{Characterize}: measure signals at the interface with an oscilloscope or multimeter and compare to your ICD specifications. (4)~\textbf{Root-cause}: determine whether the problem is hardware, software, or an interface specification error. (5)~\textbf{Fix and regression-test}: implement the fix, verify it resolves the original failure, then \textit{re-run all previously passed tests} to confirm the fix did not break anything else. \textbf{Change ONE variable at a time.} If you swap firmware and a component simultaneously, you will not know which change fixed it---and the other may have introduced a new latent defect.
\end{protip}

\begin{tipbox}
Use this sprint to finalize your test plan for the external team. Read through every test procedure as if you have never seen the project before. If any step is ambiguous, fix it now. You will not be in the room when the external team runs your tests.
\end{tipbox}

\begin{checklistbox}[Test Plan Review Checklist (Before External Testing)]
\begin{enumerate}[leftmargin=*, itemsep=1pt, topsep=0pt]
    \item Every requirement has a test procedure with a unique ID
    \item Each test has explicit, quantified pass/fail criteria (not ``works correctly'')
    \item Setup instructions include specific equipment, settings, and connections
    \item Preconditions are stated (warm-up time, calibration, initial state)
    \item Steps are numbered and unambiguous; a stranger can follow them
    \item Safety warnings appear before the step that creates the hazard
    \item Data recording format is specified (table, units, sig figs)
    \item Expected results are stated so the tester can judge pass/fail on the spot
    \item VCRM is complete with all requirements mapped to test procedures
\end{enumerate}
\end{checklistbox}

\begin{warnbox}
Do not forget to update your Bill of Materials (Part B) with actual costs, including any components you ended up not using or had to reorder. Also include a labor estimate---this is easy to overlook but is required.
\end{warnbox}

\begin{protip}
\textbf{Build a Verification Cross-Reference Matrix (VCRM) before the system demo.} The VCRM is a simple table linking each requirement to its verification method, test procedure ID, and current status (Open, Verified, Failed, Waived). By Sprint~5, you should be able to fill in most of this table. It becomes the single artifact that tells the complete project story at your final review, and it makes writing the final report dramatically easier. See MRD Chapter~5; the VCRM evolution figure shows how the matrix matures from PDR through FDR.
\end{protip}

\begin{tipbox}
\textbf{When a test fails, investigate root causes in order.} Three possible sources, checked in this sequence: (1)~Is the \textit{test} wrong? (procedure error, miscalibrated equipment, incorrect setup), (2)~Is the \textit{implementation} wrong? (the most common cause---a component, connection, or code issue), (3)~Is the \textit{requirement} wrong? (overly conservative or based on flawed assumptions). Only consider changing a requirement after ruling out the first two, and any requirement change must be formally documented with rationale. Never silently adjust a requirement to make a failed test ``pass.''
\end{tipbox}

%=============================================================================
